%%%
%%% Radical "⼞"
%%%
\section*{Radical 31: ``⼞''}\addcontentsline{toc}{section}{Radical 31: ⼞}\addcontentsline{loh}{figure}{\#\#\#\# 31: ⼞}

%%%%%%%%%% 囘 %%%%%%%%%%
\subsection*{囘}\addcontentsline{loh}{figure}{囘}

\begin{Entry}{囘}{5}{⼞}
  \begin{Phonetics}{囘}{hui2}
    \variantof{回}
  \end{Phonetics}
\end{Entry}

%%%%%%%%%% 囚 %%%%%%%%%%
\subsection*{囚}\addcontentsline{loh}{figure}{囚}

\begin{Entry}{囚}{5}{⼞}
  \begin{Phonetics}{囚}{qiu2}
    \definition[个,群,位,名,些,批]{s.}{prisioneiro; condenado}
    \definition{v.}{aprisionar}
  \end{Phonetics}
\end{Entry}

\begin{Entry}{囚犯}{5,5}{⼞,⽝}
  \begin{Phonetics}{囚犯}{qiu2fan4}[][HSK 7-9]
    \definition[名]{s.}{prisioneiro; condenado}
  \end{Phonetics}
\end{Entry}

%%%%%%%%%% 四 %%%%%%%%%%
\subsection*{四}\addcontentsline{loh}{figure}{四}

\begin{Entry}{四}{5}{⼞}
  \begin{Phonetics}{四}{si4}[][HSK 1]
    \definition*{s.}{Sobrenome: Si}
    \definition{num.}{quatro; 4}
    \definition{s.}{uma nota da escala em Gongchepu (工尺谱), correspondente ao 6 na notação musical numerada}
  \seealsoref{工尺谱}{gong1che3 pu3}
  \end{Phonetics}
\end{Entry}

\begin{Entry}{四川}{5,3}{⼞,⼮}
  \begin{Phonetics}{四川}{si4chuan1}
    \definition*{s.}{Província de Sichuan}
  \end{Phonetics}
\end{Entry}

\begin{Entry}{四处}{5,5}{⼞,⼡}
  \begin{Phonetics}{四处}{si4chu4}[][HSK 6]
    \definition{adv.}{em volta; ao redor; em todos os lugares; em todas as direções | refere"-se a um local abstrato e só pode ser colocado antes de um verbo como um adjunto adverbial}
  \seealsoref{四周}{si4zhou1}
  \synonymref{遍地}{bian4di4}
  \synonymref{处处}{chu4chu4}
  \synonymref{到处}{dao4chu4}
  \synonymref{随处}{sui2chu4}
  \synonymref{随地}{sui2di4}
  \antonymref{一面}{yi2mian4}
  \antonymref{一边}{yi4bian1}
  \end{Phonetics}
\end{Entry}

\begin{Entry}{四合院}{5,6,9}{⼞,⼝,⾩}
  \begin{Phonetics}{四合院}{si4he2yuan4}[][HSK 7-9]
    \definition[座,所,幢,套]{s.}{pátio; habitações quadrangulares; um conjunto de casas em torno de um pátio quadrado; o estilo residencial tradicional em Pequim consiste em casas nos quatro lados e um pátio no meio}
  \end{Phonetics}
\end{Entry}

\begin{Entry}{四周}{5,8}{⼞,⼝}
  \begin{Phonetics}{四周}{si4zhou1}[][HSK 5]
    \definition{s.}{ao redor; por todos os lados; a parte que circunda o centro | refere"-se a um local específico e pode ser usado não apenas como um adjunto adverbial, mas também como um ponto central}
  \seealsoref{四处}{si4chu4}
  \synonymref{周边}{zhou1bian1}
  \synonymref{周围}{zhou1wei2}
  \antonymref{中间}{zhong1jian1}
  \antonymref{中心}{zhong1xin1}
  \antonymref{中央}{zhong1yang1}
  \end{Phonetics}
\end{Entry}

\begin{Entry}{四季}{5,8}{⼞,⼦}
  \begin{Phonetics}{四季}{si4ji4}[][HSK 7-9]
    \definition{s.}{quatro estações; primavera, verão, outono e inverno, cada uma com duração de três meses}
  \end{Phonetics}
\end{Entry}

\begin{Entry}{四季分明}{5,8,4,8}{⼞,⼦,⼑,⽇}
  \begin{Phonetics}{四季分明}{si4ji4 fen1ming2}
    \definition{expr.}{as quatro estações distintas}
  \end{Phonetics}
\end{Entry}

\begin{Entry}{四季如春}{5,8,6,9}{⼞,⼦,⼥,⽇}
  \begin{Phonetics}{四季如春}{si4ji4-ru2chun1}
    \definition{expr.}{``Com clima de primavera o ano todo.''; como a primavera o ano todo; clima favorável durante todo o ano; quatro estações como a primavera}
  \antonymref{冰天雪地}{bing1tian1-xue3di4}
  \end{Phonetics}
\end{Entry}

\begin{Entry}{四面八方}{5,9,2,4}{⼞,⾯,⼋,⽅}
  \begin{Phonetics}{四面八方}{si4mian4-ba1fang1}[][HSK 7-9]
    \definition{expr.}{em todas as direções; em todos os ângulos; ao redor; de perto e de longe}
  \synonymref{大街小巷}{da4jie1-xiao3xiang4}
  \end{Phonetics}
\end{Entry}

%%%%%%%%%% 回 %%%%%%%%%%
\subsection*{回}\addcontentsline{loh}{figure}{回}

\begin{Entry}{回}{6}{⼞}
  \begin{Phonetics}{回}{hui2}[][HSK 1,2]
    \definition*{s.}{Sobrenome: Hui}
    \definition*{s.}{Etnia Hui (mulçumanos chineses)}
    \definition{clas.}{usado para coisas, ações, número de vezes |  um trecho de um conto; um capítulo de um romance em capítulos | seção ou capítulo (de um livro clássico)}
    \definition{v.}{circular; enrolar | retornar; voltar; voltar ao lugar de origem | dar meia"-volta | responder; contestar | relatar; reportar; responder}
  \seealsoref{次}{ci4}
  \seealsoref{返}{fan3}
  \synonymref{归}{gui1}
  \antonymref{来}{lai2}
  \end{Phonetics}
\end{Entry}

\begin{Entry}{回升}{6,4}{⼞,⼗}
  \begin{Phonetics}{回升}{hui2sheng1}[][HSK 7-9]
    \definition{v.}{levantar"-se novamente (após uma queda); levantar"-se | recuperar; cair e depois subir novamente}
  \antonymref{回落}{hui2luo4}
  \end{Phonetics}
\end{Entry}

\begin{Entry}{回忆}{6,4}{⼞,⼼}
  \begin{Phonetics}{回忆}{hui2yi4}[][HSK 5]
    \definition[个,段]{s.}{memória; lembrança de eventos ou experiências passadas | pode combinar"-se com adjetivos como ``doce, feliz\dots'' ou com verbos como ``desencadear'' para formar complementos | refere"-se a pensamentos mais casuais e espontâneos sobre eventos passados}
    \definition{v.}{lembrar; recordar | usado tanto na linguagem falada quanto na escrita, referindo"-se a eventos vivenciados pessoalmente}
  \seealsoref{回顾}{hui2gu4}
  \seealsoref{回想}{hui2xiang3}
  \synonymref{回首}{hui2shou3}
  \synonymref{纪念}{ji4nian4}
  \synonymref{记忆}{ji4yi4}
  \antonymref{忘却}{wang4que4}
  \antonymref{无视}{wu2shi4}
  \antonymref{展望}{zhan3wang4}
  \end{Phonetics}
\end{Entry}

\begin{Entry}{回忆录}{6,4,8}{⼞,⼼,⼹}
  \begin{Phonetics}{回忆录}{hui2yi4lu4}[][HSK 7-9]
    \definition{s.}{memórias; reminiscências; lembranças; um gênero de escrita que relata experiências pessoais ou eventos históricos familiares}
  \end{Phonetics}
\end{Entry}

\begin{Entry}{回去}{6,5}{⼞,⼛}
  \begin{Phonetics}{回去}{hui2qu5}[][HSK 1]
    \definition{v.}{retornar; voltar; estar de volta; (a partir da minha localização)}
  \synonymref{重返}{chong2fan3}
  \synonymref{返回}{fan3hui2}
  \synonymref{归来}{gui1lai2}
  \synonymref{回归}{hui2gui1}
  \synonymref{退回}{tui4hui2}
  \antonymref{出发}{chu1fa1}
  \antonymref{离开}{li2//kai1}
  \antonymref{前往}{qian2wang3}
  \antonymref{外出}{wai4chu1}
  \antonymref{走开}{zou3kai1}
  \end{Phonetics}
\end{Entry}

\begin{Entry}{回头}{6,5}{⼞,⼤}
  \begin{Phonetics}{回头}{hui2tou2}[][HSK 5]
    \definition{adv.}{mais tarde; depois de um tempo}
    \definition{conj.}{ou então; usado no início da segunda metade de uma frase para indicar o que acontecerá se você não fizer o que fez na primeira metade da frase}
    \definition{v.}{dar a meia"-volta; virar a cabeça; virar a cabeça para trás | retornar; voltar | arrepender"-se; corrigir seu caminho; reconhecer e corrigir erros}
  \synonymref{不然}{bu4ran2}
  \synonymref{否则}{fou3ze2}
  \synonymref{回顾}{hui2gu4}
  \synonymref{回家}{hui2jia1}
  \synonymref{回首}{hui2shou3}
  \synonymref{回来}{hui2lai5}
  \synonymref{醒悟}{xing3wu4}
  \synonymref{省悟}{xing3wu4}
  \end{Phonetics}
\end{Entry}

\begin{Entry}{回归}{6,5}{⼞,⼹}
  \begin{Phonetics}{回归}{hui2gui1}[][HSK 7-9]
    \definition{v.}{retornar; regredir; retornar para (local original, organização, etc.)}
  \end{Phonetics}
\end{Entry}

\begin{Entry}{回扣}{6,6}{⼞,⼿}
  \begin{Phonetics}{回扣}{hui2kou4}[][HSK 7-9]
    \definition[笔,的]{s.}{propina; desconto; comissão sobre vendas (para agente)}
  \end{Phonetics}
\end{Entry}

\begin{Entry}{回收}{6,6}{⼞,⽁}
  \begin{Phonetics}{回收}{hui2shou1}[][HSK 5]
    \definition{v.}{reciclar; reciclar itens (geralmente resíduos ou produtos antigos) para reutilização | recuperar; recolher; recuperar o que foi emitido ou demitido}
  \synonymref{收回}{shou1hui2}
  \antonymref{发布}{fa1bu4}
  \antonymref{发放}{fa1fang4}
  \antonymref{发射}{fa1she4}
  \antonymref{发送}{fa1song5}
  \end{Phonetics}
\end{Entry}

\begin{Entry}{回应}{6,7}{⼞,⼴}
  \begin{Phonetics}{回应}{hui2ying4}[][HSK 6]
    \definition{v.}{responder}
  \synonymref{答复}{da2fu5}
  \synonymref{回答}{hui2da2}
  \synonymref{回复}{hui2fu4}
  \synonymref{响应}{xiang3ying4}
  \antonymref{沉默}{chen2mo4}
  \end{Phonetics}
\end{Entry}

\begin{Entry}{回报}{6,7}{⼞,⼿}
  \begin{Phonetics}{回报}{hui2bao4}[][HSK 5]
    \definition{s.}{recompensa; pagamento; benefícios recebidos como resultado de assistência, esforço ou afeto | retornos; benefícios recebidos por meio de investimentos}
    \definition{v.}{pagar de volta; beneficiar pessoas ou organizações que os ajudaram ou cuidaram deles de alguma forma}
  \synonymref{报答}{bao4da2}
  \synonymref{报复}{bao4fu4}
  \synonymref{报告}{bao4gao4}
  \synonymref{报酬}{bao4chou5}
  \synonymref{答复}{da2fu5}
  \synonymref{反击}{fan3ji1}
  \synonymref{回馈}{hui2kui4}
  \synonymref{汇报}{hui4bao4}
  \antonymref{辜负}{gu1fu4}
  \antonymref{索取}{suo3qu3}
  \end{Phonetics}
\end{Entry}

\begin{Entry}{回来}{6,7}{⼞,⽊}
  \begin{Phonetics}{回来}{hui2lai5}[][HSK 1]
    \definition{v.}{voltar; regressar (para a minha localização) | retornar; usado após um verbo, significa ``vir ao lugar original''}
  \synonymref{重返}{chong2fan3}
  \synonymref{返回}{fan3hui2}
  \synonymref{归来}{gui1lai2}
  \synonymref{过来}{guo4//lai2}
  \synonymref{回归}{hui2gui1}
  \synonymref{回去}{hui2qu5}
  \antonymref{出去}{chu1qu5}
  \antonymref{离开}{li2//kai1}
  \antonymref{走开}{zou3kai1}
  \end{Phonetics}
\end{Entry}

\begin{Entry}{回到}{6,8}{⼞,⼑}
  \begin{Phonetics}{回到}{hui2dao4}[][HSK 1]
    \definition{v.}{retornar para; voltar e chegar (ao lugar onde estava originalmente); (após uma mudança nas circunstâncias) retornar ao estado original}
  \end{Phonetics}
\end{Entry}

\begin{Entry}{回味}{6,8}{⼞,⼝}
  \begin{Phonetics}{回味}{hui2wei4}[][HSK 7-9]
    \definition{s.}{sabor residual; recordar o sabor agradável de\dots; o gosto residual que fica na boca depois de comer}
    \definition{v.}{recordar e ponderar sobre; reviver coisas que você vivenciou ou com as quais entrou em contato}
  \end{Phonetics}
\end{Entry}

\begin{Entry}{回国}{6,8}{⼞,⼞}
  \begin{Phonetics}{回国}{hui2guo2}[][HSK 2]
    \definition{v.}{regressar ao seu país (terra natal); referindo"-se a voltar do exterior}
  \antonymref{出国}{chu1//guo2}
  \end{Phonetics}
\end{Entry}

\begin{Entry}{回信}{6,9}{⼞,⼈}
  \begin{Phonetics}{回信}{hui2//xin4}[][HSK 5]
    \definition[封]{s.}{uma carta em resposta; uma mensagem verbal em resposta}
    \definition{v.+compl.}{escrever em resposta; escrever de volta; responder uma carta; responder verbalmente uma mensagem}
  \end{Phonetics}
\end{Entry}

\begin{Entry}{回复}{6,9}{⼞,⼢}
  \begin{Phonetics}{回复}{hui2fu4}[][HSK 4]
    \definition{v.}{responder (a uma carta) | retornar ao estado normal; restaurar algo ao seu estado original}
  \synonymref{答复}{da2fu5}
  \synonymref{复原}{fu4//yuan2}
  \synonymref{还原}{huan2//yuan2}
  \synonymref{恢复}{hui1fu4}
  \synonymref{回答}{hui2da2}
  \end{Phonetics}
\end{Entry}

\begin{Entry}{回首}{6,9}{⼞,⾸}
  \begin{Phonetics}{回首}{hui2shou3}[][HSK 7-9]
    \definition{v.}{virar a cabeça; virar"-se (em volta) | olhar para trás; lembrar"-se; recordar}
  \end{Phonetics}
\end{Entry}

\begin{Entry}{回家}{6,10}{⼞,⼧}
  \begin{Phonetics}{回家}{hui2jia1}[][HSK 1]
    \definition{v.}{ir (voltar) para casa; estar em casa; voltar para casa}
  \antonymref{出门}{chu1//men2}
  \antonymref{外出}{wai4chu1}
  \end{Phonetics}
\end{Entry}

\begin{Entry}{回顾}{6,10}{⼞,⾴}
  \begin{Phonetics}{回顾}{hui2gu4}[][HSK 5]
    \definition{v.}{olhar para trás | revisar; fazer uma retrospectiva; olhar para trás, pensar no passado |o significado original é ``voltar"-se e olhar'', metaforicamente significando relembrar eventos passados | seu conteúdo não se refere apenas a experiências pessoais, mas também inclui eventos significativos da nação e da sociedade, e é frequentemente usado na linguagem falada}
  \seealsoref{回忆}{hui2yi4}
  \synonymref{回首}{hui2shou3}
  \synonymref{回头}{hui2tou2}
  \synonymref{回想}{hui2xiang3}
  \synonymref{记忆}{ji4yi4}
  \synonymref{总结}{zong3jie2}
  \antonymref{展望}{zhan3wang4}
  \end{Phonetics}
\end{Entry}

\begin{Entry}{回旋}{6,11}{⼞,⽅}
  \begin{Phonetics}{回旋}{hui2xuan2}
    \definition{v.}{dar voltas em torno de; dar voltas; girar; rodar | (espaço para) manobrar}
  \end{Phonetics}
\end{Entry}

\begin{Entry}{回答}{6,12}{⼞,⽵}
  \begin{Phonetics}{回答}{hui2da2}[][HSK 1]
    \definition[个]{s.}{resposta}
    \definition{v.}{responder; explicar a questão; expressar opinião sobre a solicitação}
  \seealsoref{答案}{da2'an4}
  \seealsoref{答复}{da2fu5}
  \synonymref{回复}{hui2fu4}
  \synonymref{回应}{hui2ying4}
  \synonymref{解答}{jie3da2}
  \antonymref{提问}{ti2wen4}
  \antonymref{询问}{xun2wen4}
  \end{Phonetics}
\end{Entry}

\begin{Entry}{回落}{6,12}{⼞,⾋}
  \begin{Phonetics}{回落}{hui2luo4}[][HSK 7-9]
    \definition{v.}{(níveis de água, preços, etc.) cair após uma subida; diminuir | recuar | retornar ao nível baixo após uma subida (no nível da água, preço etc.)}
  \antonymref{回升}{hui2sheng1}
  \end{Phonetics}
\end{Entry}

\begin{Entry}{回馈}{6,12}{⼞,⾶}
  \begin{Phonetics}{回馈}{hui2kui4}[][HSK 7-9]
    \definition{v.}{retribuir; recompensar | dar \emph{feedback}; fornecer \emph{feedback}; dar retorno; dar parecer}
  \end{Phonetics}
\end{Entry}

\begin{Entry}{回想}{6,13}{⼞,⼼}
  \begin{Phonetics}{回想}{hui2xiang3}[][HSK 7-9]
    \definition{v.}{recordar; pensar de novo; pensar (no passado) | enfatiza a recordação de memórias profundas, reencenações mentais conscientes e tem um significado mais forte em comparação com 回忆}
  \seealsoref{回忆}{hui2yi4}
  \synonymref{回顾}{hui2gu4}
  \synonymref{回首}{hui2shou3}
  \synonymref{记忆}{ji4yi4}
  \synonymref{追溯}{zhui1su4}
  \antonymref{展望}{zhan3wang4}
  \end{Phonetics}
\end{Entry}

\begin{Entry}{回避}{6,16}{⼞,⾌}
  \begin{Phonetics}{回避}{hui2bi4}[][HSK 5]
    \definition{v.}{fugir (de um problema); em direito, refere"-se especificamente à não participação nos procedimentos de um caso de um oficial de justiça, etc., que tenha interesse no caso ou nas partes do caso | esquivar"-se; evadir"-se; evitar (encontrar alguém)}
  \synonymref{躲避}{duo3bi4}
  \synonymref{逃避}{tao2bi4}
  \antonymref{接近}{jie1jin4}
  \antonymref{面对}{mian4dui4}
  \antonymref{面临}{mian4lin2}
  \antonymref{应对}{ying4dui4}
  \antonymref{正视}{zheng4shi4}
  \antonymref{重视}{zhong4shi4}
  \end{Phonetics}
\end{Entry}

%%%%%%%%%% 因 %%%%%%%%%%
\subsection*{因}\addcontentsline{loh}{figure}{因}

\begin{Entry}{因}{6}{⼞}
  \begin{Phonetics}{因}{yin1}[][HSK 6]
    \definition*{s.}{Sobrenome: Yin}
    \definition{conj.}{porque; orações de conexão, indicando relações de causa e efeito}
    \definition{prep.}{com base em; à luz de; de acordo com; a introdução da ação comportamental equivale a 按照 ou 根据}
    \definition{s.}{causa; motivo; condições em que algo ocorre ou causa um determinado resultado}
    \definition{v.}{seguir; continuar; fazer como sempre fez | estar em conformidade com; estar de acordo com; depender; contar com}
  \seealsoref{按照}{an4zhao4}
  \seealsoref{根据}{gen1ju4}
  \synonymref{据}{ju4}
  \synonymref{凭}{ping2}
  \synonymref{袭}{xi2}
  \synonymref{沿}{yan2}
  \antonymref{果}{guo3}
  \end{Phonetics}
\end{Entry}

\begin{Entry}{因人而异}{6,2,6,6}{⼞,⼈,⽽,⼶}
  \begin{Phonetics}{因人而异}{yin1ren2'er2yi4}[][HSK 7-9]
    \definition{expr.}{``Isso varia de pessoa para pessoa.''; varia de pessoa para pessoa; pessoas diferentes se comportam de maneira diferente; métodos diferentes devem ser usados dependendo do objeto}
  \end{Phonetics}
\end{Entry}

\begin{Entry}{因为}{6,4}{⼞,⼂}
  \begin{Phonetics}{因为}{yin1wei5}[][HSK 2]
    \definition{conj.}{porque; indica o motivo e a frase seguinte indica o resultado | usado para expressar o propósito de uma ação | geralmente usado na primeira metade de uma frase para expressar razão ou causa | 因为 e 因而 não podem aparecer simultaneamente em uma frase e não podem substituir um ao outro}
    \definition{prep.}{por causa de; por conta de; indica razão ou justificativa}
  \seealsoref{为了}{wei4le5}
  \seealsoref{因而}{yin1'er2}
  \synonymref{鉴于}{jian4yu2}
  \synonymref{由于}{you2yu2}
  \antonymref{所以}{suo3yi3}
  \end{Phonetics}
\end{Entry}

\begin{Entry}{因为……所以……}{6,4,8,4}{⼞,⼂,⼾,⼈}
  \begin{Phonetics}{因为……所以……}{yin1wei5 suo3yi3}
    \definition{conj.}{porque\dots portanto\dots}
  \end{Phonetics}
\end{Entry}

\begin{Entry}{因此}{6,6}{⼞,⽌}
  \begin{Phonetics}{因此}{yin1ci3}[][HSK 3]
    \definition{conj.}{assim; portanto; consequentemente}
  \seealsoref{所以}{suo3yi3}
  \seealsoref{于是}{yu2shi4}
  \synonymref{从而}{cong2'er2}
  \synonymref{为此}{wei4ci3}
  \synonymref{因而}{yin1'er2}
  \synonymref{由此}{you2ci3}
  \end{Phonetics}
\end{Entry}

\begin{Entry}{因此就}{6,6,12}{⼞,⽌,⼪}
  \begin{Phonetics}{因此就}{yin1ci3 jiu4}
    \definition{conj.}{portanto}
  \end{Phonetics}
\end{Entry}

\begin{Entry}{因而}{6,6}{⼞,⽽}
  \begin{Phonetics}{因而}{yin1'er2}[][HSK 5]
    \definition{conj.}{assim; como resultado; com o resultado que; conecta frases, indicando relação de causa e efeito | usado na segunda metade para expressar resultado, equivalente a 所以 e 因此 | 因为 e 因而 não podem aparecer simultaneamente em uma frase e não podem substituir um ao outro}
  \seealsoref{所以}{suo3yi3}
  \seealsoref{因此}{yin1ci3}
  \seealsoref{因为}{yin1wei5}
  \synonymref{从而}{cong2'er2}
  \synonymref{为此}{wei4ci3}
  \synonymref{由此}{you2ci3}
  \synonymref{于是}{yu2shi4}
  \end{Phonetics}
\end{Entry}

\begin{Entry}{因素}{6,10}{⼞,⽷}
  \begin{Phonetics}{因素}{yin1su4}[][HSK 6]
    \definition[个,种]{s.}{fator; elemento; os componentes que constituem a essência das coisas | fator; as razões ou condições que determinam o sucesso ou o fracasso de algo}
  \synonymref{成分}{cheng2fen5}
  \synonymref{要素}{yao4su4}
  \end{Phonetics}
\end{Entry}

%%%%%%%%%% 团 %%%%%%%%%%
\subsection*{团}\addcontentsline{loh}{figure}{团}

\begin{Entry}{团}{6}{⼞}
  \begin{Phonetics}{团}{tuan2}[][HSK 3]
    \definition*{s.}{Liga da Juventude Comunista da China; Liga}
    \definition{adj.}{redondo; circular}
    \definition{clas.}{usado para algo em forma de bola}
    \definition[个]{s.}{bolinho de massa; comida em forma de bola feita de arroz ou farinha | algo em forma de bola | grupo; corpo; sociedade; organização; um grupo envolvido em um determinado trabalho ou atividade | regimento; unidade organizacional militar, geralmente abaixo do nível de divisão e acima do nível de batalhão}
    \definition{v.}{enrolar algo para formar uma bola; rolar | reunir; unir; conglomerar}
  \synonymref{凑}{cou4}
  \synonymref{会}{hui4}
  \synonymref{集}{ji2}
  \synonymref{聚}{ju4}
  \synonymref{圆}{yuan2}
  \antonymref{分}{fen1}
  \antonymref{散}{san3}
  \end{Phonetics}
\end{Entry}

\begin{Entry}{团长}{6,4}{⼞,⾧}
  \begin{Phonetics}{团长}{tuan2zhang3}[][HSK 5]
    \definition[位,名]{s.}{comandante do regimento | chefe (ou presidente) de uma delegação, trupe, etc. | líder de uma delegação}
  \end{Phonetics}
\end{Entry}

\begin{Entry}{团队}{6,4}{⼞,⾩}
  \begin{Phonetics}{团队}{tuan2dui4}[][HSK 6]
    \definition[个,支,种]{s.}{equipe; time; grupo; um grupo de alguma natureza}
  \synonymref{队伍}{dui4wu5}
  \synonymref{集体}{ji2ti3}
  \synonymref{群体}{qun2ti3}
  \synonymref{任务}{ren4wu5}
  \synonymref{社团}{she4tuan2}
  \synonymref{团体}{tuan2ti3}
  \synonymref{小组}{xiao3zu3}
  \synonymref{组织}{zu3zhi1}
  \antonymref{个人}{ge4ren2}
  \antonymref{孤立}{gu1li4}
  \end{Phonetics}
\end{Entry}

\begin{Entry}{团伙}{6,6}{⼞,⼈}
  \begin{Phonetics}{团伙}{tuan2huo3}[][HSK 7-9]
    \definition[个]{s.}{gangue; bando; panelinha | gangue (criminosa) | cúmplice; comparsa; membro de gangue}
  \end{Phonetics}
\end{Entry}

\begin{Entry}{团体}{6,7}{⼞,⼈}
  \begin{Phonetics}{团体}{tuan2ti3}[][HSK 3]
    \definition[种,个]{s.}{equipe; grupo; organização; um grupo de pessoas com objetivos e interesses comuns}
  \seealsoref{集体}{ji2ti3}
  \synonymref{集团}{ji2tuan2}
  \synonymref{全体}{quan2ti3}
  \synonymref{社团}{she4tuan2}
  \synonymref{团队}{tuan2dui4}
  \synonymref{协会}{xie2hui4}
  \synonymref{整体}{zheng3ti3}
  \antonymref{个人}{ge4ren2}
  \antonymref{个体}{ge4ti3}
  \end{Phonetics}
\end{Entry}

\begin{Entry}{团员}{6,7}{⼞,⼝}
  \begin{Phonetics}{团员}{tuan2yuan2}[][HSK 7-9]
    \definition[名,位,个]{s.}{membro; membro de um grupo específico | membro da Liga; membro da Liga da Juventude Comunista da China; refere"-se especificamente aos membros da Liga da Juventude Comunista da China}
  \synonymref{会员}{hui4yuan2}
  \synonymref{团聚}{tuan2ju4}
  \end{Phonetics}
\end{Entry}

\begin{Entry}{团结}{6,9}{⼞,⽷}
  \begin{Phonetics}{团结}{tuan2jie2}[][HSK 3]
    \definition{adj.}{unido; amigável; harmonioso; relação harmoniosa e coexistência harmoniosa}
    \definition{v.}{unir; reunir}
  \synonymref{合并}{he2bing4}
  \synonymref{合作}{he2zuo4}
  \synonymref{互助}{hu4zhu4}
  \synonymref{结合}{jie2he2}
  \synonymref{联合}{lian2he2}
  \synonymref{统一}{tong3yi1}
  \synonymref{协作}{xie2zuo4}
  \antonymref{斗争}{dou4zheng1}
  \antonymref{分开}{fen1//kai1}
  \antonymref{分离}{fen1li2}
  \antonymref{分裂}{fen1lie4}
  \antonymref{分散}{fen1san4}
  \antonymref{瓦解}{wa3jie3}
  \end{Phonetics}
\end{Entry}

\begin{Entry}{团圆}{6,10}{⼞,⼞}
  \begin{Phonetics}{团圆}{tuan2yuan2}[][HSK 7-9]
    \definition{adj.}{redondo; indica que a forma é circular}
    \definition{v.}{reunir; reencontrar membros da família após um período de separação}
  \synonymref{团聚}{tuan2ju4}
  \antonymref{分离}{fen1li2}
  \antonymref{离别}{li2bie2}
  \end{Phonetics}
\end{Entry}

\begin{Entry}{团聚}{6,14}{⼞,⽿}
  \begin{Phonetics}{团聚}{tuan2ju4}[][HSK 7-9]
    \definition{v.}{reunir; reencontrar familiares ou amigos muito próximos após a separação | unir; reunir; congregar}
  \synonymref{重逢}{chong2feng2}
  \synonymref{欢聚}{huan1ju4}
  \synonymref{聚会}{ju4hui4}
  \synonymref{团圆}{tuan2yuan2}
  \synonymref{团员}{tuan2yuan2}
  \antonymref{分开}{fen1//kai1}
  \antonymref{分别}{fen1bie2}
  \antonymref{分离}{fen1li2}
  \antonymref{分散}{fen1san4}
  \antonymref{离别}{li2bie2}
  \end{Phonetics}
\end{Entry}

%%%%%%%%%% 园 %%%%%%%%%%
\subsection*{园}\addcontentsline{loh}{figure}{园}

\begin{Entry}{园}{7}{⼞}
  \begin{Phonetics}{园}{yuan2}[][HSK 6]
    \definition*{s.}{Sobrenome: Yuan}
    \definition[个]{s.}{jardim; terreno; plantação; terra para cultivar plantas | local para recreação pública; locais para passeios turísticos e entretenimento | área para fins especiais | uma área de terra para o cultivo de plantas; um lugar onde vegetais, flores, frutas e árvores são cultivados}
  \end{Phonetics}
\end{Entry}

\begin{Entry}{园地}{7,6}{⼞,⼟}
  \begin{Phonetics}{园地}{yuan2di4}[][HSK 6]
    \definition{s.}{jardim; campo | Figurativo: campo; escopo}
  \synonymref{场地}{chang3di4}
  \synonymref{殿堂}{dian4tang2}
  \synonymref{农场}{nong2chang3}
  \synonymref{天地}{tian1di4}
  \synonymref{田园}{tian2yuan2}
  \synonymref{舞台}{wu3tai2}
  \synonymref{园林}{yuan2lin2}
  \antonymref{空地}{kong1di4}
  \antonymref{空地}{kong4di4}
  \end{Phonetics}
\end{Entry}

\begin{Entry}{园林}{7,8}{⼞,⽊}
  \begin{Phonetics}{园林}{yuan2lin2}[][HSK 5]
    \definition[处,座,个]{s.}{parque; jardim; área paisagística com plantas e árvores para as pessoas apreciarem e descansarem.}
  \synonymref{花园}{hua1yuan2}
  \synonymref{绿地}{lv4di4}
  \end{Phonetics}
\end{Entry}

%%%%%%%%%% 囯 %%%%%%%%%%
\subsection*{囯}\addcontentsline{loh}{figure}{囯}

\begin{Entry}{囯}{7}{⼞}
  \begin{Phonetics}{囯}{guo2}
    \definition*{s.}{Sobrenome: Guo}
    \definition{adj.}{do estado; nacional | do nosso país; Chinês | do país}
    \definition{s.}{país; nação; estado | o melhor da nação | o melhor; o mais bonito do país}
    \variantof{国}
  \end{Phonetics}
\end{Entry}

%%%%%%%%%% 困 %%%%%%%%%%
\subsection*{困}\addcontentsline{loh}{figure}{困}

\begin{Entry}{困}{7}{⼞}
  \begin{Phonetics}{困}{kun4}[][HSK 3]
    \definition{adj.}{cansado; exausto; fatigado | difícil; complicado; difícil e penoso; pobre e miserável | sonolento; com sono; cansado, com vontade de dormir}
    \definition{v.}{ficar encalhado; estar em apuros; preso em dificuldades e sofrimentos ou limitado por circunstâncias e condições que não pode escapar | cercar; envolver; imobilizar; controlar dentro de um determinado limite | dormir}
  \end{Phonetics}
\end{Entry}

\begin{Entry}{困厄}{7,4}{⼞,⼚}
  \begin{Phonetics}{困厄}{kun4'e4}
    \definition{adj.}{(viver) em privação | (situação) difícil | em águas profundas}
  \synonymref{窘迫}{jiong3po4}
  \synonymref{困境}{kun4jing4}
  \synonymref{逆境}{ni4jing4}
  \end{Phonetics}
\end{Entry}

\begin{Entry}{困扰}{7,7}{⼞,⼿}
  \begin{Phonetics}{困扰}{kun4rao3}[][HSK 5]
    \definition{v.}{perturbar; deixar perplexo; perseguir}
  \synonymref{苦恼}{ku3nao3}
  \synonymref{困惑}{kun4huo4}
  \synonymref{扰乱}{rao3luan4}
  \synonymref{疑惑}{yi2huo4}
  \synonymref{忧虑}{you1lv4}
  \end{Phonetics}
\end{Entry}

\begin{Entry}{困苦}{7,8}{⼞,⾋}
  \begin{Phonetics}{困苦}{kun4ku3}
    \definition{adj.}{(viver) em privação}
    \definition{s.}{dificuldade; privação; tribulação; sofrimento; miséria; angústia; pobreza}
  \seealsoref{困厄}{kun4'e4}
  \seealsoref{痛苦}{tong4ku3}
  \synonymref{艰巨}{jian1ju4}
  \synonymref{艰苦}{jian1ku3}
  \synonymref{艰难}{jian1nan2}
  \synonymref{困难}{kun4nan5}
  \synonymref{难过}{nan2guo4}
  \synonymref{贫困}{pin2kun4}
  \synonymref{贫穷}{pin2qiong2}
  \synonymref{辛苦}{xin1ku3}
  \antonymref{舒适}{shu1shi4}
  \antonymref{甜蜜}{tian2mi4}
  \antonymref{幸福}{xing4fu2}
  \end{Phonetics}
\end{Entry}

\begin{Entry}{困难}{7,10}{⼞,⾫}
  \begin{Phonetics}{困难}{kun4nan5}[][HSK 3]
    \definition{adj.}{dificuldades financeiras; circunstâncias difíceis | complicado; complexo; difícil; árduo; a situação é complexa e há muitos obstáculos}
    \definition[种]{s.}{dificuldade; situação difícil; problemas ou situações difíceis de resolver no trabalho e na vida}
  \seealsoref{艰难}{jian1nan2}
  \synonymref{艰苦}{jian1ku3}
  \synonymref{麻烦}{ma2fan5}
  \synonymref{难处}{nan2chu4}
  \synonymref{难度}{nan2du4}
  \synonymref{难关}{nan2guan1}
  \synonymref{难处}{nan2chu5}
  \antonymref{便利}{bian4li4}
  \antonymref{富裕}{fu4yu4}
  \antonymref{简单}{jian3dan1}
  \antonymref{简易}{jian3yi4}
  \antonymref{轻易}{qing1yi5}
  \antonymref{容易}{rong2yi4}
  \antonymref{顺利}{shun4li4}
  \end{Phonetics}
\end{Entry}

\begin{Entry}{困惑}{7,12}{⼞,⼼}
  \begin{Phonetics}{困惑}{kun4huo4}[][HSK 7-9]
    \definition{adj.}{em um quebra"-cabeça; me sinto confuso e sem saber o que fazer}
    \definition{v.}{sentir"-se perplexo; ser intrigante o porquê, e eu não saber o que fazer}
  \end{Phonetics}
\end{Entry}

\begin{Entry}{困境}{7,14}{⼞,⼟}
  \begin{Phonetics}{困境}{kun4jing4}[][HSK 7-9]
    \definition{s.}{dilema; situação difícil; apuro}
  \end{Phonetics}
\end{Entry}

%%%%%%%%%% 围 %%%%%%%%%%
\subsection*{围}\addcontentsline{loh}{figure}{围}

\begin{Entry}{围}{7}{⼞}
  \begin{Phonetics}{围}{wei2}[][HSK 3]
    \definition*{s.}{Sobrenome: Wei}
    \definition{clas.}{o comprimento das duas mãos com os polegares e os dedos indicadores juntos ou dos dois braços juntos}
    \definition{s.}{em volta de tudo; ao redor}
    \definition{v.}{cercar; rodear; circundar; encurralar; cercar tudo, impedindo a passagem entre o interior e o exterior | envolver; contornar | usado na linguagem falada}
  \seealsoref{包围}{bao1wei2}
  \synonymref{环}{huan2}
  \synonymref{圈}{quan1}
  \synonymref{绕}{rao4}
  \synonymref{周}{zhou1}
  \antonymref{开}{kai1}
  \antonymref{通}{tong1}
  \end{Phonetics}
\end{Entry}

\begin{Entry}{围巾}{7,3}{⼞,⼱}
  \begin{Phonetics}{围巾}{wei2jin1}[][HSK 4]
    \definition[条]{s.}{lenço; cachecol; echarpe; gravata; tiras longas de malha ou tecido usadas ao redor do pescoço para aquecimento, proteção do colarinho ou decoração}
  \end{Phonetics}
\end{Entry}

\begin{Entry}{围绕}{7,9}{⼞,⽷}
  \begin{Phonetics}{围绕}{wei2rao4}[][HSK 5]
    \definition{v.}{girar; circundar; dar voltas; girar em torno de algo; cercar | concentrar"-se em; centrar"-se em; centrar"-se em uma questão ou evento (para realizar atividades)}
  \synonymref{环绕}{huan2rao4}
  \end{Phonetics}
\end{Entry}

\begin{Entry}{围棋}{7,12}{⼞,⽊}
  \begin{Phonetics}{围棋}{wei2qi2}
    \definition*{s.}{Go; Weiqi; um jogo de tabuleiro com peças pretas e brancas; o jogador com mais peças ocupando mais espaços vence}
  \end{Phonetics}
\end{Entry}

\begin{Entry}{围墙}{7,14}{⼞,⼟}
  \begin{Phonetics}{围墙}{wei2qiang2}[][HSK 7-9]
    \definition[堵,道,面]{s.}{recinto; parede envolvente; construído em torno de uma casa, jardim, pátio, etc.; uma parede que serve de barreira}
  \end{Phonetics}
\end{Entry}

%%%%%%%%%% 固 %%%%%%%%%%
\subsection*{固}\addcontentsline{loh}{figure}{固}

\begin{Entry}{固}{8}{⼞}
  \begin{Phonetics}{固}{gu4}
    \definition*{s.}{Sobrenome: Gu}
    \definition{adj.}{sólido; firme; forte | duro; sólido | mal informado; superficial; ignorante}
    \definition{adv.}{firmemente; resolutamente | originalmente; em primeiro lugar | certamente; reconhecidamente; seguramente}
    \definition{conj.}{usado da mesma forma que 固然}
    \definition{v.}{solidificar; consolidar; fortalecer | defender; proteger}
  \seealsoref{固然}{gu4ran2}
  \end{Phonetics}
\end{Entry}

\begin{Entry}{固执}{8,6}{⼞,⼿}
  \begin{Phonetics}{固执}{gu4zhi5}[][HSK 7-9]
    \definition{adj.}{obstinado; teimoso; mantém suas próprias opiniões e não quer mudá"-las, mesmo que estejam erradas}
  \end{Phonetics}
\end{Entry}

\begin{Entry}{固定}{8,8}{⼞,⼧}
  \begin{Phonetics}{固定}{gu4ding4}[][HSK 4]
    \definition{adj.}{fixo; regular; inalterado ou imóvel}
    \definition{v.}{consertar; tornar fixo, não mover novamente; colocar as coisas em ordem, não mudá"-las novamente}
  \synonymref{静止}{jing4zhi3}
  \synonymref{确定}{que4ding4}
  \synonymref{锁定}{suo3ding4}
  \synonymref{停止}{ting2zhi3}
  \synonymref{稳定}{wen3ding4}
  \synonymref{稳固}{wen3gu4}
  \antonymref{变换}{bian4huan4}
  \antonymref{动摇}{dong4yao2}
  \antonymref{流动}{liu2dong4}
  \antonymref{调节}{tiao2jie2}
  \antonymref{摇摆}{yao2bai3}
  \antonymref{摇晃}{yao2huang4}
  \antonymref{转化}{zhuan3hua4}
  \antonymref{转移}{zhuan3yi2}
  \end{Phonetics}
\end{Entry}

\begin{Entry}{固然}{8,12}{⼞,⽕}
  \begin{Phonetics}{固然}{gu4ran2}[][HSK 7-9]
    \definition{conj.}{usado para introduzir uma cláusula adversativa admitindo primeiro um certo fato; quando usado na primeira metade de uma frase, a segunda metade geralmente tem 可是 ou 但是 para ecoá"-lo, indicando que o fato A é reconhecido, mas o fato B não se torna inválido por causa do fato A | admitir um fato sem negar outro; indica o reconhecimento de um fato, levando a uma transição no texto seguinte; indica o reconhecimento do fato A e não nega o fato B}
  \seealsoref{但是}{dan4shi4}
  \seealsoref{可是}{ke3shi4}
  \end{Phonetics}
\end{Entry}

%%%%%%%%%% 国 %%%%%%%%%%
\subsection*{国}\addcontentsline{loh}{figure}{国}

\begin{Entry}{国}{8}{⼞}
  \begin{Phonetics}{国}{guo2}[][HSK 1]
    \definition*{s.}{Sobrenome: Guo}
    \definition{adj.}{nacional; do estado; representante do país | o melhor de um país}
    \definition[个,片]{s.}{estado; nação; país}
  \seealsoref{国家}{guo2jia1}
  \end{Phonetics}
\end{Entry}

\begin{Entry}{国人}{8,2}{⼞,⼈}
  \begin{Phonetics}{国人}{guo2ren2}
    \definition{s.}{compatriota; conterrâneo; pessoas do mesmo país}
  \end{Phonetics}
\end{Entry}

\begin{Entry}{国土}{8,3}{⼞,⼟}
  \begin{Phonetics}{国土}{guo2tu3}[][HSK 7-9]
    \definition{s.}{terra; território; território nacional}
  \end{Phonetics}
\end{Entry}

\begin{Entry}{国内}{8,4}{⼞,⼌}
  \begin{Phonetics}{国内}{guo2nei4}[][HSK 3]
    \definition{s.}{interno (a um país); doméstico; lar; dentro de um determinado país}
  \synonymref{境内}{jing4nei4}
  \antonymref{国际}{guo2ji4}
  \antonymref{国外}{guo2wai4}
  \antonymref{海外}{hai3wai4}
  \antonymref{境外}{jing4wai4}
  \end{Phonetics}
\end{Entry}

\begin{Entry}{国王}{8,4}{⼞,⽟}
  \begin{Phonetics}{国王}{guo2wang2}[][HSK 6]
    \definition[位,名,个,些]{s.}{rei; soberanos; o governante supremo de algumas monarquias antigas; nos tempos modernos, refere"-se ao chefe de estado de algumas monarquias}
  \end{Phonetics}
\end{Entry}

\begin{Entry}{国外}{8,5}{⼞,⼣}
  \begin{Phonetics}{国外}{guo2wai4}[][HSK 1]
    \definition{adj.}{externo; no exterior; fora do país; outros lugares fora do país; geralmente chamados de exterior;  exterior não é o mesmo que estrangeiro}
  \synonymref{海外}{hai3wai4}
  \antonymref{国内}{guo2nei4}
  \end{Phonetics}
\end{Entry}

\begin{Entry}{国民}{8,5}{⼞,⽒}
  \begin{Phonetics}{国民}{guo2min2}[][HSK 5]
    \definition{adj.}{nacional}
    \definition[个]{s.}{membro de uma nação; povo de uma nação}
  \synonymref{公民}{gong1min2}
  \synonymref{民众}{min2zhong4}
  \synonymref{群众}{qun2zhong4}
  \synonymref{人民}{ren2min2}
  \end{Phonetics}
\end{Entry}

\begin{Entry}{国产}{8,6}{⼞,⼇}
  \begin{Phonetics}{国产}{guo2chan3}[][HSK 6]
    \definition{adj.}{doméstico; feito na China; produzido internamente, especificamente na China}
  \antonymref{进口}{jin4//kou3}
  \end{Phonetics}
\end{Entry}

\begin{Entry}{国会}{8,6}{⼞,⼈}
  \begin{Phonetics}{国会}{guo2hui4}[][HSK 6]
    \definition{s.}{parlamento; congresso}
  \synonymref{议会}{yi4hui4}
  \end{Phonetics}
\end{Entry}

\begin{Entry}{国庆}{8,6}{⼞,⼴}
  \begin{Phonetics}{国庆}{guo2qing4}[][HSK 3]
    \definition*{s.}{Dia Nacional, o dia em que um país comemora sua independência ou fundação}
  \end{Phonetics}
\end{Entry}

\begin{Entry}{国庆节}{8,6,5}{⼞,⼴,⾋}
  \begin{Phonetics}{国庆节}{guo2qing4jie2}
    \definition*[个]{s.}{Dia Nacional, refere"-se ao dia em que um país celebra sua independência ou fundação; na China, o Dia Nacional é comemorado em 1º de outubro de cada ano para celebrar a fundação da República Popular da China}
  \end{Phonetics}
\end{Entry}

\begin{Entry}{国有}{8,6}{⼞,⽉}
  \begin{Phonetics}{国有}{guo2you3}[][HSK 7-9]
    \definition{v.}{pertencer ao estado; ser nacionalizado}
  \end{Phonetics}
\end{Entry}

\begin{Entry}{国防}{8,6}{⼞,⾩}
  \begin{Phonetics}{国防}{guo2fang2}[][HSK 7-9]
    \definition{s.}{defesa nacional; as instalações humanas, materiais e militares que um país possui para defender sua soberania territorial e impedir invasões estrangeiras}
  \end{Phonetics}
\end{Entry}

\begin{Entry}{国际}{8,7}{⼞,⾩}
  \begin{Phonetics}{国际}{guo2ji4}[][HSK 2]
    \definition{adj.}{internacional; entre países; entre nações}
    \definition{s.}{internacional; o mundo; entre nações; entre países de todo o mundo}
  \synonymref{全球}{quan2qiu2}
  \synonymref{世界}{shi4jie4}
  \antonymref{地方}{di4fang1}
  \antonymref{国内}{guo2nei4}
  \antonymref{局部}{ju2bu4}
  \antonymref{区域}{qu1yu4}
  \end{Phonetics}
\end{Entry}

\begin{Entry}{国际儿童节}{8,7,2,12,5}{⼞,⾩,⼉,⽴,⾋}
  \begin{Phonetics}{国际儿童节}{guo2ji4 er2tong2jie2}
    \definition*{s.}{Dia Internacional das Crianças (1 de junho)}
  \end{Phonetics}
\end{Entry}

\begin{Entry}{国际妇女节}{8,7,6,3,5}{⼞,⾩,⼥,⼥,⾋}
  \begin{Phonetics}{国际妇女节}{guo2ji4fu4nv3jie2}
    \definition*{s.}{Dia Internacional das Mulheres (8 de março)}
  \end{Phonetics}
\end{Entry}

\begin{Entry}{国际劳动节}{8,7,7,6,5}{⼞,⾩,⼒,⼒,⾋}
  \begin{Phonetics}{国际劳动节}{guo2ji4 lao2dong4jie2}
    \definition*{s.}{Dia Internacional dos Trabalhadores (1 de maio)}
  \end{Phonetics}
\end{Entry}

\begin{Entry}{国学}{8,8}{⼞,⼦}
  \begin{Phonetics}{国学}{guo2xue2}[][HSK 7-9]
    \definition*{s.}{Arcaico: O Colégio Imperial}
    \definition{s.}{estudos da cultura clássica chinesa (história, filosofia, literatura, língua, etc.) | cultura nacional chinesa | estudos da antiga civilização chinesa}
  \end{Phonetics}
\end{Entry}

\begin{Entry}{国宝}{8,8}{⼞,⼧}
  \begin{Phonetics}{国宝}{guo2bao3}[][HSK 7-9]
    \definition[件]{s.}{tesouro nacional}
  \end{Phonetics}
\end{Entry}

\begin{Entry}{国画}{8,8}{⼞,⽥}
  \begin{Phonetics}{国画}{guo2hua4}[][HSK 7-9]
    \definition[幅,张,卷]{s.}{pintura tradicional chinesa | arte chinesa | pintura nacional}
  \end{Phonetics}
\end{Entry}

\begin{Entry}{国语}{8,9}{⼞,⾔}
  \begin{Phonetics}{国语}{guo2yu3}
    \definition*{s.}{Língua Chinesa (Mandarim), enfatizando sua natureza nacional}
  \antonymref{方言}{fang1yan2}
  \antonymref{外语}{wai4yu3}
  \end{Phonetics}
\end{Entry}

\begin{Entry}{国家}{8,10}{⼞,⼧}
  \begin{Phonetics}{国家}{guo2jia1}[][HSK 1]
    \definition[个]{s.}{país; estado; nação; um lugar reconhecido internacionalmente e com soberania independente, incluindo as pessoas e as instituições administrativas desse lugar}
  \seealsoref{国}{guo2}
  \synonymref{国土}{guo2tu3}
  \end{Phonetics}
\end{Entry}

\begin{Entry}{国宾馆}{8,10,11}{⼞,⼧,⾷}
  \begin{Phonetics}{国宾馆}{guo2bin1guan3}
    \definition{s.}{pousada estatal | casa de hóspedes estatal}
  \end{Phonetics}
\end{Entry}

\begin{Entry}{国情}{8,11}{⼞,⼼}
  \begin{Phonetics}{国情}{guo2qing2}[][HSK 7-9]
    \definition{s.}{condição (ou estado) do país; condições nacionais; as condições e características básicas da natureza social, política, economia, cultura etc. de um país também se referem especificamente às condições e características básicas de um país em um determinado período de tempo}
  \end{Phonetics}
\end{Entry}

\begin{Entry}{国旗}{8,14}{⼞,⽅}
  \begin{Phonetics}{国旗}{guo2qi2}[][HSK 6]
    \definition[面]{s.}{bandeira (de um país)}
  \end{Phonetics}
\end{Entry}

\begin{Entry}{国歌}{8,14}{⼞,⽋}
  \begin{Phonetics}{国歌}{guo2ge1}[][HSK 6]
    \definition[首,支]{s.}{hino nacional; o hino nacional da China, oficialmente designado pelo estado como a música que representa o país, é ``Marcha dos Voluntários''}
  \end{Phonetics}
\end{Entry}

\begin{Entry}{国徽}{8,17}{⼞,⼻}
  \begin{Phonetics}{国徽}{guo2hui1}[][HSK 7-9]
    \definition{s.}{emblema nacional; o emblema nacional da China, oficialmente designado pelo estado para representar o país, apresenta a Praça da Paz Celestial sob o céu brilhante de cinco estrelas, cercada por espigas de grãos e engrenagens}
  \end{Phonetics}
\end{Entry}

\begin{Entry}{国籍}{8,20}{⼞,⽵}
  \begin{Phonetics}{国籍}{guo2ji2}[][HSK 5]
    \definition[个]{s.}{nacionalidade; cidadania; refere"-se à identidade de um indivíduo como pertencente a um Estado | identidade nacional (de um avião, navio, etc.)}
  \end{Phonetics}
\end{Entry}

%%%%%%%%%% 图 %%%%%%%%%%
\subsection*{图}\addcontentsline{loh}{figure}{图}

\begin{Entry}{图}{8}{⼞}
  \begin{Phonetics}{图}{tu2}[][HSK 3]
    \definition*{s.}{Sobrenome: Tu}
    \definition[张,幅]{s.}{mapa; gráfico; imagem; desenho | plano; esquema; tentativa}
    \definition{v.}{procurar; perseguir; esperar obter| desenhar; retratar; pintar | imaginar; planejar; pensar; maquinar}
  \synonymref{画}{hua4}
  \synonymref{绘}{hui4}
  \synonymref{谋}{mou2}
  \synonymref{贪}{tan1}
  \synonymref{像}{xiang4}
  \antonymref{舍}{she3}
  \antonymref{文}{wen2}
  \antonymref{予}{yu3}
  \antonymref{字}{zi4}
  \end{Phonetics}
\end{Entry}

\begin{Entry}{图书}{8,4}{⼞,⼄}
  \begin{Phonetics}{图书}{tu2shu1}[][HSK 6]
    \definition{s.}{livros; um termo geral para publicações como livros e álbuns de imagens}[这些图书都可以借阅。===Esses livros estão disponíveis para empréstimo.]
  \end{Phonetics}
\end{Entry}

\begin{Entry}{图书馆}{8,4,11}{⼞,⼄,⾷}
  \begin{Phonetics}{图书馆}{tu2shu1guan3}[][HSK 1]
    \definition[家,座,个]{s.}{biblioteca; instituição que coleta, organiza e armazena livros e materiais para leitura e consulta}
  \end{Phonetics}
\end{Entry}

\begin{Entry}{图片}{8,4}{⼞,⽚}
  \begin{Phonetics}{图片}{tu2pian4}[][HSK 2]
    \definition[张,幅]{s.}{imagem; fotografia; um termo geral para imagens, fotografias, decalques, etc. usados para ilustrar algo}
  \synonymref{图画}{tu2hua4}
  \synonymref{相片}{xiang4pian4}
  \synonymref{照片}{zhao4pian4}
  \end{Phonetics}
\end{Entry}

\begin{Entry}{图形}{8,7}{⼞,⼺}
  \begin{Phonetics}{图形}{tu2xing2}[][HSK 7-9]
    \definition[种,个]{s.}{figura; gráfico | sinal; carta; desenho; diagrama}
  \synonymref{图案}{tu2'an4}
  \end{Phonetics}
\end{Entry}

\begin{Entry}{图纸}{8,7}{⼞,⽷}
  \begin{Phonetics}{图纸}{tu2zhi3}[][HSK 7-9]
    \definition[张]{s.}{desenho; planta; folha de desenho; \emph{blueprint}; um documento técnico que utiliza desenhos e texto para descrever a estrutura, forma, dimensões e outros requisitos de estruturas de engenharia, máquinas, equipamentos, etc.}
  \end{Phonetics}
\end{Entry}

\begin{Entry}{图画}{8,8}{⼞,⽥}
  \begin{Phonetics}{图画}{tu2hua4}[][HSK 3]
    \definition[幅,张,套]{s.}{desenho; imagem; pintura}
  \synonymref{绘画}{hui4hua4}
  \synonymref{图案}{tu2'an4}
  \end{Phonetics}
\end{Entry}

\begin{Entry}{图表}{8,8}{⼞,⾐}
  \begin{Phonetics}{图表}{tu2biao3}[][HSK 7-9]
    \definition[张,个]{s.}{carta; diagrama; gráfico | figura; gráfico; diagrama; pictograma; pictografia; cronograma; folha; tabela; termo geral para diagramas e tabelas que representam diversas situações e indicam vários números, como diagramas esquemáticos e tabelas estatísticas}
  \end{Phonetics}
\end{Entry}

\begin{Entry}{图标}{8,9}{⼞,⽊}
  \begin{Phonetics}{图标}{tu2biao1}
    \definition{s.}{ícone (informática)}[这个图标代表设置。===Este ícone representa as configurações.]
  \end{Phonetics}
\end{Entry}

\begin{Entry}{图案}{8,10}{⼞,⽊}
  \begin{Phonetics}{图案}{tu2'an4}[][HSK 4]
    \definition[种,个]{s.}{padrão; desenho; padrões e gráficos usados para decoração de edifícios, tecidos, artes e artesanato, etc.}
  \synonymref{图形}{tu2xing2}
  \end{Phonetics}
\end{Entry}

\begin{Entry}{图像}{8,13}{⼞,⼈}
  \begin{Phonetics}{图像}{tu2xiang4}[][HSK 7-9]
    \definition{s.}{imagem; figura; imagens desenhadas, fotografadas ou impressas}
  \end{Phonetics}
\end{Entry}

%%%%%%%%%% 圆 %%%%%%%%%%
\subsection*{圆}\addcontentsline{loh}{figure}{圆}

\begin{Entry}{圆}{10}{⼞}
  \begin{Phonetics}{圆}{yuan2}[][HSK 4]
    \definition*{s.}{Sobrenome: Yuan}
    \definition{adj.}{redondo; circular; esférico; arredondado | diplomático; satisfatório}
    \definition[个,轮]{s.}{círculo; circunferência | uma moeda de valor e peso fixos}
    \definition{v.}{tornar plausível; justificar; tornar completo; completar}
  \antonymref{扁}{bian3}
  \antonymref{方}{fang1}
  \end{Phonetics}
\end{Entry}

\begin{Entry}{圆形}{10,7}{⼞,⼺}
  \begin{Phonetics}{圆形}{yuan2xing2}[][HSK 7-9]
    \definition{adj.}{circular; redondo}
    \definition{s.}{círculo; a figura formada pela circunferência de um círculo em um plano | com formato de bola ou hemisfério}
  \end{Phonetics}
\end{Entry}

\begin{Entry}{圆珠笔}{10,10,10}{⼞,⽟,⽵}
  \begin{Phonetics}{圆珠笔}{yuan2zhu1bi3}[][HSK 6]
    \definition[支,枝]{s.}{caneta esferográfica}
  \end{Phonetics}
\end{Entry}

\begin{Entry}{圆满}{10,13}{⼞,⽔}
  \begin{Phonetics}{圆满}{yuan2man3}[][HSK 4]
    \definition{adj.}{perfeito; satisfatório; sem defeitos}
  \synonymref{美满}{mei3man3}
  \synonymref{完备}{wan2bei4}
  \synonymref{完满}{wan2man3}
  \synonymref{完美}{wan2mei3}
  \synonymref{完善}{wan2shan4}
  \antonymref{悲惨}{bei1can3}
  \antonymref{不足}{bu4zu2}
  \antonymref{欠缺}{qian4que1}
  \antonymref{缺陷}{que1xian4}
  \end{Phonetics}
\end{Entry}

%%%%%%%%%% 圈 %%%%%%%%%%
\subsection*{圈}\addcontentsline{loh}{figure}{圈}

\begin{Entry}{圈}{11}{⼞}
  \begin{Phonetics}{圈}{juan1}
    \definition{v.}{prender aves e animais de criação | Coloquial: prender os criminosos; colocar na cadeia, prisão | confinar; encarcerar}
  \end{Phonetics}
  \begin{Phonetics}{圈}{juan4}[][HSK 7-9]
    \definition*{s.}{Sobrenome: Juan}
    \definition{s.}{curral; local onde o gado ou as aves são mantidos, geralmente cercado ou murado, alguns com galpões}
  \end{Phonetics}
  \begin{Phonetics}{圈}{quan1}[][HSK 4]
    \definition[个]{s.}{anel; círculo; refere"-se a algo em forma de anel | domínio; grupo; escopo; círculo(s)}
    \definition{v.}{cercar; rodear; circundar | marcar com um círculo}
  \end{Phonetics}
\end{Entry}

\begin{Entry}{圈子}{11,3}{⼞,⼦}
  \begin{Phonetics}{圈子}{quan1zi5}[][HSK 7-9]
    \definition[个]{s.}{anel; círculo; uma figura plana, redonda e oca; um objeto em forma de anel | círculo; panelinha; grupo; âmbito; refere"-se ao âmbito das atividades humanas ou à área de um grupo}
  \end{Phonetics}
\end{Entry}

\begin{Entry}{圈套}{11,10}{⼞,⼤}
  \begin{Phonetics}{圈套}{quan1tao4}[][HSK 7-9]
    \definition[个]{s.}{malha; armadilha; laço; esquemas para enganar pessoas}
  \end{Phonetics}
\end{Entry}

\begin{Entry}{圈粉}{11,10}{⼞,⽶}
  \begin{Phonetics}{圈粉}{quan1fen3}
    \definition{s.}{Neologismo, Coloquial: ganhar alguém como fã, obter novos fãs}
  \end{Phonetics}
\end{Entry}

%%%%% EOF %%%%%

